% 13-vra.tex — Clause 13: Value range constraints
\chapter{Value range constraints}
\label{sec:13-vra}
\index{value range analysis (VRA)}
\index{VRA}

\section{Overview}
\label{sec:vra.overview}

\pnum
Value Range Analysis (VRA) is a static analysis that tracks integer value
ranges through a program using interval arithmetic.  It is used to:
\begin{itemize}
  \item eliminate array bounds checks (\S\,\ref{sec:types.array});
  \item prove division-by-zero absence (\S\,\ref{sec:expr.arith});
  \item verify parameter and return constraints (\S\,\ref{sec:vra.constraints});
  \item prove in-guard safety (\S\,\ref{sec:expr.in}).
\end{itemize}

\pnum
VRA is decidable and polynomial in program size.  No SMT solver is required.
All analyses performed by a conforming implementation shall terminate.

\section{Range representation}
\label{sec:vra.range}
\index{range representation}

\pnum
A value range is an interval $[\ell,\, u]$ where $\ell \leq u$.  The
sentinels $-\infty$ and $+\infty$ represent the absence of a lower or upper
bound, respectively.  A range may additionally be \textit{unknown} (no
information available), in which case the interval is $(-\infty, +\infty)$.

\section{Range propagation rules}
\label{sec:vra.propagation}
\index{range propagation}

\pnum
The following table defines the result range for binary arithmetic
operations on intervals $[a,b]$ and $[c,d]$.  All arithmetic is
saturating: overflow wraps to $\pm\infty$.

\begin{longtable}{@{}lll@{}}
\toprule
\textbf{Operation} & \textbf{Result range} & \textbf{Notes} \\
\midrule
\endfirsthead
\toprule
\textbf{Operation} & \textbf{Result range} & \textbf{Notes} \\
\midrule
\endhead
\bottomrule
\endlastfoot
$[a,b] + [c,d]$ & $[a+c,\; b+d]$ & saturating \\
$[a,b] - [c,d]$ & $[a-d,\; b-c]$ & saturating \\
$[a,b] \times [c,d]$ & $[\min(ac,ad,bc,bd),\;\max(ac,ad,bc,bd)]$ & saturating \\
$[a,b] \div [c,d]$, $c > 0$ & $[a \div d,\; b \div c]$ & integer truncation \\
$[a,b] \div [c,d]$, $0 \in [c,d]$ & \diagcode{E015} & division by zero \\
$[a,b] \mathbin{\%} [c,d]$, $a \geq 0$, $c > 0$ & $[0,\; d-1]$ & \\
$[a,b] \mathbin{\%} [c,d]$, otherwise & unknown & conservative \\
\end{longtable}

\section{Conditional refinement}
\label{sec:vra.refinement}
\index{conditional refinement}

\pnum
After an \lain{if x < N} test, the VRA refines the range of \lain{x}
in the then-branch to $[\ell,\, N-1]$ and in the else-branch to $[N,\, u]$
(where $[\ell, u]$ is the range before the test).

\pnum
At the join point after an \lain{if}/\lain{else}, the range is the union
of the ranges from both branches.

\pnum
When conditions are chained with \lain{and}, refinements from the
left operand are active when evaluating the right operand
(\S\,\ref{sec:expr.in}).

\section{In-guard semantics}
\label{sec:vra.in-guard}
\index{in guard!VRA semantics}

\pnum
An in-guard \lain{idx in arr} (\S\,\ref{sec:expr.in}) introduces the
constraint $0 \leq \mathit{idx} < \mathit{arr.len}$ into the VRA's
range environment.  Inside the guarded body, \lain{arr[idx]} is accepted
without further proof.

\pnum
When \lain{idx in arr} is the condition of a \lain{while decreasing M}
loop, the VRA infers that $\mathit{arr.len} - \mathit{idx} \geq 0$ (a
valid non-negativity proof for the measure $\mathit{arr.len} - \mathit{idx}$).

\section{Constraint annotations}
\label{sec:vra.constraints}
\index{constraint annotation}

\pnum
A constraint annotation of the form \lain{T op bound} on a parameter or
return type declares an invariant that the VRA shall enforce:

\begin{itemize}
  \item \textbf{Parameter constraint}: at every call site, the VRA shall
        prove that the argument satisfies the constraint.
  \item \textbf{Return constraint}: at every \lain{return} statement, the
        VRA shall prove that the return expression satisfies the constraint.
\end{itemize}

\begin{constraint}
  A function call where the argument range cannot be proven to satisfy
  a parameter constraint is \illformed{}~(\diagcode{E012}).
\end{constraint}

\begin{constraint}
  A \lain{return} expression whose range cannot be proven to satisfy the
  declared return constraint is \illformed{}~(\diagcode{E086}) --- the same
  boundary mechanism (\S\,\ref{sec:vra.boundary}) that checks assignments and
  narrowing conversions.  (A \emph{parameter}-constraint violation at a call
  site is the distinct code \diagcode{E012}, above.)
\end{constraint}

\section{Interprocedural range propagation}
\label{sec:vra.interprocedural}
\index{interprocedural VRA}

\pnum
When a call expression appears in a context where the VRA needs a range,
it uses the callee's declared return constraint (if any) as a conservative
approximation.

\begin{example}
\begin{laincode}
func nonzero() int >= 1 { return 42 }

func divide(a int) int {
    return a / nonzero()   // divisor range known >= 1: no [E015]
}
\end{laincode}
\end{example}

\pnum
When the callee has no declared return constraint, the call expression has
an unknown range.

\begin{note}
  Interprocedural VRA is conservative: only declared constraints are used,
  not body analysis.  This preserves decidability and ensures that changing
  a function's body cannot silently affect its callers' safety proofs.
\end{note}

\section{Limits and known gaps}
\label{sec:vra.limits}

\pnum
The following limitations are known and accepted:

\begin{enumerate}
  \item VRA does not track general integer value ranges of struct fields: a
        field's value range is unknown (a field read through a pointer always
        is).  Two narrower field properties \emph{are} tracked --- per-field
        \emph{initialization} (definite assignment, \S\,\ref{sec:own.diag}) and
        a struct field's \emph{slice length} within an \lain{and}-scope (to
        prove indexing of a field slice).
  \item VRA does not track ranges across loop-carried dependences with
        non-linear measures.
  \item Integer overflow in VRA arithmetic uses saturating semantics; the
        actual Lain program uses wrapping arithmetic.  Programs whose
        correctness depends on integer overflow are not analyzed correctly
        by VRA.
  \item Floating-point values are not tracked by VRA.
\end{enumerate}


\section{Loop body affine recap}
\label{sec:vra.affine}
\index{affine loop body}

\pnum
For \lain{for V in start..end} loops whose body contains assignments of
the form \lain{x = x + c} or \lain{x = x - c} with constant step~\(c\),
the implementation computes the post-loop range of~\(x\) precisely:

\begin{align*}
  x_{\text{post}} \in [x_{\text{init}} + c \cdot n_{\min},\; x_{\text{init}} + c \cdot n_{\max}]
\end{align*}

where \(n_{\min}\), \(n_{\max}\) is the iteration count derived from the
loop's range. This avoids the conservative widening that would otherwise
mark \(x\) as unknown after the loop.

\begin{example}
\begin{laincode}
var x = 0
for i in 0..10 {
    x = x + 1
}
// Post-loop: x is exactly 10 (range [10, 10]).
\end{laincode}
\end{example}

\pnum
Patterns not matching the affine form (mixed updates in branches,
non-constant step, multiple variables updated together) fall back to
the conservative widening (unknown post-loop range).

\section{Overflow-at-boundary check}
\label{sec:vra.boundary}
\index{overflow check}

\pnum
The Phase 5 boundary check applies at every integer-type boundary
where a value crosses into a narrower or different-domain target.
The check sites are:
\begin{itemize}
  \item Variable declaration (\lain{var x T = expr}).
  \item Assignment to identifier, field, or indexed element
        (\lain{x = expr}, \lain{p.f = expr}, \lain{arr[i] = expr}).
  \item Return value vs the function's declared return type.
  \item Call argument vs the parameter's declared type.
  \item Struct constructor field initialization.
\end{itemize}

\pnum
At each site, if the source value's VRA range exceeds the target type's
natural range, the program is \illformed{}~(\diagcode{E086}). The
diagnostic offers these resolutions: widen the target; choose an explicit
overflow policy via a cast tier or operator --- checked (\lain{as?},
\lain{+?}), wrapping (\lain{as\%}, \lain{+\%}), or saturating (\lain{as|},
\lain{+|}); or tighten input constraints so VRA can prove safety
(\S\,\ref{sec:types.overflow}).

\pnum
The check is skipped:
\begin{itemize}
  \item Inside an \lain{unsafe} block;
  \item When the source range is unbounded (the analysis cannot rule out
        overflow without further information; the user must annotate);
  \item For integer-literal sources (literals are polymorphic across
        \lain{iN}/\lain{uN} and trusted to fit when written explicitly).
\end{itemize}

\pnum
The wrapping/saturating operators (\tok{+\%}, \tok{*\%}, \tok{+|}, \tok{*|})
produce a result whose range is automatically clamped to the LHS type's
natural range, so they pass the boundary check unconditionally.

\begin{example}
\begin{laincode}
var a u8 = 200
var b u8 = 100
// var bad u8 = a + b      // [E086]: range [300, 300] exceeds u8
var w u8 = a +% b           // OK: wrapping clamps to [0, 255]
var s u8 = a +| b           // OK: saturating clamps to [0, 255]
\end{laincode}
\end{example}

\section{Sub-slice bounds}
\label{sec:vra.subslice}
\index{sub-slice}
\index{slicing}

\pnum
A sub-slice expression \lain{arr[a..b]} produces a slice covering the
half-open range $[a, b)$ of \lain{arr}.  The inferred type of the
result is \lain{T[b-a]}, a sized slice
(\S\,\ref{sec:types.sized-slice}) whose length is statically known.
When both endpoints are compile-time integer literals, the compiler
verifies $a \leq b$; if the check fails the program is
\illformed{}~(\diagcode{E087}).

\pnum
When the endpoints are variables, each is range-checked against the
container length (\diagcode{E085}, the standard bounds check). A
runtime guarantee on the ordering of variable endpoints is not
emitted at compile time; the programmer must establish it through
refinement constraints or by limiting the range expression to
literal bounds.

\begin{example}
\begin{laincode}
var arr i32[10] = [0,1,2,3,4,5,6,7,8,9]
var ok  = arr[2..5]    // type i32[3], half-open, length 3
// var bad = arr[5..2] // [E087]: start 5 > end 2
var empty = arr[3..3]  // OK: empty slice of length 0
\end{laincode}
\end{example}

\section{Sized slice call-site verification}
\label{sec:vra.sized-slice-callsite}
\index{sized slice!call-site check}
\index{E087}

\pnum
When a function parameter is declared with a sized slice type
(\S\,\ref{sec:types.sized-slice}), the compiler verifies at every
call site that the supplied argument satisfies the parameter's
length constraint.  A proved violation is \illformed{}~(\diagcode{E087}).
The check is conservative: it fires only when VRA can \emph{prove}
the constraint is not met; it is silent when the constraint cannot
be determined statically.

\pnum
The cases checked are:

\begin{itemize}
  \item \lain{T[>= k]} or \lain{T[> k]} with integer literal \lain{k}:
        fires if the argument length's VRA upper bound is less than
        the required minimum (\lain{>= k} requires $\mathit{len} \geq k$;
        \lain{> k} requires $\mathit{len} \geq k+1$).
  \item \lain{T[k]} (exact length, literal \lain{k}): fires if the
        argument length's VRA range does not contain \lain{k}.
  \item \lain{T[ref.len]} (cross-parameter equality): fires if both
        sides have a concrete singleton VRA range and they differ.
\end{itemize}

\begin{example}
\begin{laincode}
func pair_sum(arr i32[>= 2]) i32 { return arr[0] + arr[1] }

proc main() i32 {
    var one i32[1] = [42]
    // pair_sum(one)        // [E087]: len 1 < required 2
    var two i32[2] = [1, 2]
    return pair_sum(two)    // OK: len 2 satisfies >= 2
}
\end{laincode}
\end{example}

\pnum
\textbf{Sized slice VRA entry.}
At function entry, for each parameter \lain{name T[expr]} the compiler
seeds the range table with the constraint derived from the size
expression:

\begin{itemize}
  \item \lain{T[k]} (literal): \texttt{\_\_len\_name} $\in [k, k]$.
  \item \lain{T[>= k]} (literal): \texttt{\_\_len\_name} $\in [k, +\infty)$.
  \item \lain{T[> k]} (literal): \texttt{\_\_len\_name} $\in [k{+}1, +\infty)$.
  \item \lain{T[a.len]} (member reference): difference constraint
        \texttt{\_\_len\_name} $=$ \texttt{\_\_len\_a}.
  \item \lain{T[a.len + b.len]} (arithmetic): constraints
        \texttt{\_\_len\_a} $\leq$ \texttt{\_\_len\_name} and
        \texttt{\_\_len\_b} $\leq$ \texttt{\_\_len\_name}.
\end{itemize}

These constraints are available to the VRA throughout the function body,
enabling the Omega Test (\S\,\ref{sec:vra.omega}) to discharge complex
indexing patterns.

\section{VRA analysis levels}
\label{sec:vra.levels}
\index{VRA!analysis levels}

\pnum
Lain's VRA applies four escalating analysis levels at each bounds-check
site.  The levels are tried in order; if one succeeds (the check is proved
safe), the remaining levels are skipped.

\begin{description}[leftmargin=2em]
  \item[L1 — Interval arithmetic.]
    Each variable is represented by an interval $[\ell, u]$.
    The check passes if the interval of the index falls within
    $[0, \mathit{arr.len} - 1]$ directly.

  \item[L2 — Difference constraints.]
    A system of difference constraints of the form
    $x - y \leq c$ is maintained.  The check passes if the constraint
    system can derive the required bound in one step.

  \item[L3 — One-step bridge (transitivity).]
    Two-hop transitivity: if $v_1 - X \leq a$ and $X - v_2 \leq b$
    are both known, the checker derives $v_1 - v_2 \leq a + b$ without
    running the full Fourier-Motzkin algorithm.  Handles patterns such
    as \lain{i < src.len - 1} via the chain
    $i < \mathit{src.len}$ and $\mathit{src.len} - 1 < \mathit{src.len}$.

  \item[L4 — Fourier-Motzkin linear arithmetic (Omega Test).]
    Full linear arithmetic elimination over all VRA variables.
    See \S\,\ref{sec:vra.omega}.
\end{description}

\pnum
All four levels are polynomial time.  No SMT solver is used.

\section{Omega Test — Fourier-Motzkin prover (VRA level 4)}
\label{sec:vra.omega}
\index{Omega Test}
\index{Fourier-Motzkin}
\index{VRA!level 4}

\pnum
When L1–L3 fail to discharge a bounds check, the implementation invokes
the \textit{Omega Test}: a self-contained Fourier-Motzkin linear
arithmetic prover (roughly 410 lines in \texttt{src/sema/omega.h}, no external
dependencies).

\pnum
\textbf{Proof method (proof by contradiction).}
To prove $\mathit{index} < \mathit{bound}$:

\begin{enumerate}
  \item Decompose both expressions into linear forms
        $\sum_i c_i \cdot x_i + k$ over named variables.
        References to \lain{x.len} are mapped to the synthetic VRA
        variable \texttt{\_\_len\_x}.
  \item Form the negated query: $\mathit{bound} - \mathit{index} \leq 0$.
  \item Extract all VRA range and difference constraints that mention
        the variables appearing in the query.
  \item Run Fourier-Motzkin variable elimination.  Each step crosses
        upper-bound inequalities with lower-bound inequalities to
        produce a system with one fewer variable.
  \item If the resulting constant system contains an inequality
        $0 \leq c$ with $c < 0$, the system is UNSAT: the negation
        is impossible, so $\mathit{index} < \mathit{bound}$ always holds.
\end{enumerate}

\pnum
The prover also supports non-negativity queries: to prove
$\mathit{expr} \geq 0$, it shows that $\mathit{expr} \leq -1$, combined
with VRA constraints, is UNSAT.

\pnum
\textbf{Example patterns proved by the Omega Test.}

\begin{example}
\begin{laincode}
// concat: out[j + a.len] with j < b.len and out.len == a.len + b.len
// FM eliminates j; the cross-pair (j < b.len, j >= b.len) gives 0 <= -1 (UNSAT).
func concat(a i32[], b i32[], out i32[a.len + b.len]) {
    for j in 0..b.len {
        out[j + a.len] = b[j]   // proved safe by Omega
    }
}

// reverse: out[n - i - 1] with n = arr.len and i < n
// FM eliminates n and __len_arr; reduces to (i <= -1) AND (i >= 0) (UNSAT).
func reverse(arr i32[], out i32[arr.len]) {
    var n = arr.len
    for i in 0..n {
        out[n - i - 1] = arr[i]   // proved safe by Omega
    }
}
\end{laincode}
\end{example}

\pnum
\textbf{Equality propagation for local length aliases.}
When the body contains \lain{var n = x.len}, the compiler registers the
equality $n - \mathtt{\_\_len\_x} = 0$ (expressed as two difference
constraints) so the Omega Test can cancel \lain{n} and
\texttt{\_\_len\_x} in the same linear form.

\pnum
\textbf{Conservatism.}
If the FM buffer overflows (more than 80 inequalities in the working
system), the prover gives up and returns \textit{unknown}; the check
then falls through to E085 as if L4 had not fired.  This bound is
unreachable for the small systems (2--6 variables, 5--15 inequalities)
typical of slice manipulation code.
